\documentclass[10pt,a4paper]{amsart}
\usepackage[margin=19mm]{geometry}
\usepackage{amsmath,amssymb,mathtools}
\usepackage[scr=rsfs,cal=euler]{mathalfa}
\usepackage{xfrac}
\usepackage[unicode,hidelinks,bookmarks=false]{hyperref}
\newcommand{\ie}{i.e.\ }
\newcommand{\cf}{cf.\ }
\newcommand{\resp}{respectively\ }
\newcommand{\Subsection}{\subsection}
\providecommand{\nobreakdash}{\nobreak\hbox{-}}
\setlength{\emergencystretch}{2em}
\multlinegap0pt
\usepackage{amssymb}
\usepackage{braket}
\usepackage{xcolor}
\usepackage{tikz}
\newtheorem{proposition}{Proposition}
\newcommand{\dd}{\mathrm{d}}
\title{Power convexity of the torsion function on horo-convex domains in hyperbolic space}
\author{Wei Zhang, Qi Zhou}
\date{Source: arXiv:2609.02516v1 (CC BY 4.0)}
\begin{document}
\maketitle

\noindent\emph{Source and licence.} Power convexity of the torsion function on horo-convex domains in hyperbolic space. Version arXiv:2609.02516v1 (CC BY 4.0), distributed by arXiv under the Creative Commons Attribution 4.0 International licence, http://creativecommons.org/licenses/by/4.0/, as recorded in the arXiv licence metadata for this exact version and shown on the abstract page https://arxiv.org/abs/2609.02516v1. Full licence notice: \texttt{licenses/arxiv-2609.02516-CC-BY-4.0.txt}.\par\smallskip

\noindent\textit{Editorial scope.} Counted unit: the article's proposition case (source0.tex lines 180--186). The article's complete original proof of it is delivered with it (source0.tex lines 187--205, one \texttt{proof} environment of 19 lines). No mathematics is written, repaired or re-derived: the statement and the proof are the article's own lines, in the article's own notation, and no retained line is substituted or re-flowed.  Retained uncounted as the source-local context the proof needs: lines 159--161; lines 175--177.  Nothing was omitted from any retained span, so the package carries no omission record.  No retained line contains a citation, so the wrapper carries no bibliography and no bibliography entry.  No retained line contains a figure environment, an image-inclusion command or drawing code, so no image is delivered and the package declares no asset dependency.  Editorial formatting only, in addition: an AMS \texttt{amsart} preamble that restates the article's own declarations and macros used by the retained lines, namely the environments \texttt{proposition} on separate counters, together with the standard packages those lines need.  The retained environments print in this document as Proposition 1 in order of appearance, whereas the article prints the same items with the numbering of its own full document.  The complete native source of the article as distributed in the arXiv e-print for this exact version is bundled unchanged as \texttt{sources/209208/source0.tex}.
\begin{equation}\label{Equality-radial Hessian}
\nabla^2u=u''\dd r^2+u'\coth r(g-\dd r^2).
\end{equation}

\begin{equation}\label{Equation-torsion-radial}
\varphi'=-(n-1)\coth r\varphi-\frac{1+\varphi^2}{v_B}.
\end{equation}

\begin{proposition}\label{Proposition-ball case}
Define
\begin{equation}
\rho_0(n)=\mathrm{arccoth}\left(\frac{1+\sqrt{5+4/(n-1)}}{2}\right).
\end{equation}
If $0<R<\rho_0(n)$, then we have $$\nabla^2v_B-|\nabla v_B|g>0\quad \mathrm{in}\ B_R.$$
\end{proposition}

\begin{proof}
By \eqref{Equality-radial Hessian}, the radial eigenvalue of $\nabla^2v_B-|\nabla v_B|g$ is $\varphi'-\varphi$, while its tangential eigenvalue is $(\coth r-1)\varphi$ with multiplicity $n-1$. Since $\varphi>0$ and $\coth r>1$, it remains to prove that $$\varphi'-\varphi>0.$$

Let $h=\varphi'-\varphi$. At the center, since $\varphi(0)=0$, taking the limit $r\rightarrow 0$ in \eqref{Equation-torsion-radial} and applying L'H\^{o}pital's rule yield $$h(0)=\varphi'(0)=-\frac{1}{nv_B(0)}>0.$$ Arguing by contradiction, suppose that $h$ vanishes somewhere in $(0,R)$. Denoting its first zero by $r_1$, one has $$h(r)>0\quad\mathrm{for}\ 0\leq r<r_1,\quad h(r_1)=0,\quad h'(r_1)\leq 0,$$ while $h(r_1)=0$ means that $\varphi'(r_1)=\varphi(r_1)$, and \eqref{Equation-torsion-radial} at $r_1$ takes the form
\begin{equation}\label{Equation-torsion-radial-r_1}
(1+(n-1)\coth r_1)\varphi=-\frac{1+\varphi^2}{v_B}.
\end{equation}
Differentiating \eqref{Equation-torsion-radial} gives $$\varphi''=(n-1)\mathrm{csch}^2r\varphi-(n-1)\coth r\varphi'-\frac{2}{v_B}\varphi\varphi'+\frac{\varphi(1+\varphi^2)}{v_B^2},$$ which, together with \eqref{Equation-torsion-radial-r_1} and $\varphi'(r_1)=\varphi(r_1)$, yields at $r_1$
\begin{align*}
h'
=&~\varphi''-\varphi'\\
=&~(n-1)\mathrm{csch}^2r_1\varphi-(n-1)\coth r_1\varphi-\frac{2}{v_B}\varphi^2+\frac{\varphi(1+\varphi^2)}{v_B^2}-\varphi\\
=&~\varphi\bigg((n-1)\mathrm{csch}^2r_1-(1+(n-1)\coth r_1)\\
&\hspace{0.3cm}+\frac{2(1+(n-1)\coth r_1)\varphi^2}{1+\varphi^2}+\frac{(1+(n-1)\coth r_1)^2\varphi^2}{1+\varphi^2}\bigg)\\
\geq&~\varphi\left((n-1)\mathrm{csch}^2r_1-1-(n-1)\coth r_1\right).
\end{align*}
Moreover, in view of the identity $\mathrm{csch}^2r_1=\coth^2r_1-1$, we have $$h'\geq(n-1)\varphi\left(\coth^2r_1-\coth r_1-\frac{n}{n-1}\right).$$ The right-hand side is positive whenever $$\coth r_1>\frac{1+\sqrt{5+4/(n-1)}}{2}.$$ Since $$\rho_0(n)=\mathrm{arccoth}\left(\frac{1+\sqrt{5+4/(n-1)}}{2}\right),$$ for $R<\rho_0(n)$, it follows from $r_1<R<\rho_0(n)$ and the strict decrease of $\coth r$ that $$\coth r_1>\coth\rho_0(n).$$
This yields $h'(r_1)>0$, whereas $h'(r_1)\leq 0$, a contradiction. Hence $h>0$ on $[0,R)$, and the proposition follows.
\end{proof}

\end{document}
